% \begin{frame}
%     \frametitle{Ответы на замечания ведущей организации НИИ~<<Рога~и~копыта>>}
%     \begin{itemize}
%         \item Замечание -- ответ
%         \item Замечание -- ответ
%         \item Замечание -- ответ
%         \item Замечание -- ответ
%         \item Замечание -- ответ
%     \end{itemize}
% \end{frame}

% \begin{frame}
%     \frametitle{Ответы на замечания оф. оппонента Иванова\,И.\,И}
%     \begin{itemize}
%         \item Замечание -- ответ
%         \item Замечание -- ответ
%         \item Замечание -- ответ
%         \item Замечание -- ответ
%         \item Замечание -- ответ
%     \end{itemize}
% \end{frame}

% \begin{frame}
%     \frametitle{Ответы на замечания Петрова\,П.\,П}
%     \begin{itemize}
%         \item Замечание -- ответ
%         \item Замечание -- ответ
%         \item Замечание -- ответ
%         \item Замечание -- ответ
%         \item Замечание -- ответ
%     \end{itemize}
% \end{frame}

% TODO: удалить после итоговой аттестации
\begin{frame}
    \frametitle{Ответы на замечания рецензента Кулакова К. А.}
    \scriptsize
    \textbf{1. Независимость операций и~бригады настройщиков.}

    Модель предполагает, что каждая машина обслуживается собственным оператором, что~соответствует типовой организации ЦБП.
    Время переналадки (перестановка ножей, смена марки) включено в~длительность операции на~конкретной машине.
    Учёт ограниченного числа бригад как~разделяемого ресурса трансформирует задачу в~задачу календарного планирования с~ресурсными ограничениями и~является направлением дальнейших исследований.

    \vspace{0.15cm}
    \textbf{2. Обоснование выбора генетического алгоритма.}

    Задача является NP-трудной (гл.~2). Глобальный план $G=(\overline{\overline{X}},\overline{\overline{Y}},\overline{\overline{Z}})$ состоит из~трёх взаимозависимых компонент~--- пространство состояний для~динамического программирования экспоненциально, а~свойство оптимальной подструктуры не~выполняется: изменение $\overline{\overline{Y}}$ требует полного пересчёта $\overline{\overline{X}}$ и~$\overline{\overline{Z}}$.
    Генетический алгоритм естественно оперирует в~пространстве глобальных планов через операторы мутации и~отбора, сохраняя допустимость решений.

    \vspace{0.15cm}
    \textbf{3. Сравнительный анализ с~работами Воронова и~Урбана.}

    Постановки указанных работ не~включают ряд элементов настоящей модели (табл.~4 диссертации): многослойные заказы, БРС, динамику складов, минимизацию максимального числа раскроев~--- прямое сравнение результатов невозможно из-за~различий в~формулировках задач, входных данных и~критериях.
    Сравнение с~полным перебором на~\emph{той~же} модели~--- стандартный подход верификации эвристик на~малых экземплярах в~комбинаторной оптимизации.
\end{frame}